\chapter{Conclusion}
\label{chap:conclusion}

This research asked whether heterogeneous literature data can support reliable machine-learning prediction of material ignition and flame spread rate for fire-safe spacecraft design, especially when prediction is required for an experimental source that the model has not previously seen. The precise answer is conditional. The database supports strong interpolation among rows from represented papers, but only moderate and substantially less stable generalisation to held-out papers. It is therefore suitable as the basis of a traceable screening and research-prioritisation tool, but it does not yet support universal prediction, causal claims, or replacement of qualification and certification testing.

For FSR regression, 2,531 usable observations from 66 canonical papers were evaluated with frozen random-row and paper-grouped holdouts. The selected random-forest baseline achieved \(R^2=0.879\pm0.018\), RMSE \(=0.00626\pm0.00047~\mathrm{m\,s^{-1}}\), and MAE \(=0.00210\pm0.00017~\mathrm{m\,s^{-1}}\) under random-row interpolation. When complete papers were withheld, the same model achieved \(R^2=0.458\pm0.222\), RMSE \(=0.01221\pm0.00412~\mathrm{m\,s^{-1}}\), and MAE \(=0.00482\pm0.00114~\mathrm{m\,s^{-1}}\). For ignition classification, 4,487 deduplicated observations from 87 papers were evaluated. The paper-square-root-weighted XGBoost interpolation champion achieved ROC--AUC \(=0.908\pm0.010\) and PR--AUC \(=0.969\pm0.004\); the unweighted XGBoost grouped champion achieved ROC--AUC \(=0.718\pm0.044\), PR--AUC \(=0.877\pm0.030\), and Brier score \(=0.179\pm0.017\).

The gap between these protocols is the primary result. Random cross-validation answers whether a model can recover nearby conditions in campaigns it has partly observed. Paper-grouped validation asks whether it can transfer across a publication boundary that bundles apparatus, laboratory practice, sampled regime, and reporting convention. The latter is the more relevant internal evidence for a new campaign, although it is not a substitute for external validation. Quoting only the random-split scores would materially overstate readiness.

The model comparison also answers the algorithmic part of the question. Tree ensembles are well matched to these mixed, sparse, nonlinear tabular data. XGBoost led ignition prediction, while random forest was selected for FSR; KNN, SVM/SVR, decision-tree, and MLP candidates supplied useful baselines but did not displace the grouped champions. This is consistent with tabular-learning literature, while also demonstrating that XGBoost is not universally best \cite{shwartzZiv2022tabular}. The wire-specific \(R^2>0.9\) findings of Rivera/Jose and colleagues remain compatible with the present results because they concern random cross-validation within a narrower electrical-wire database, not paper-disjoint transfer \cite{riveraWireML}.

Feature attributions identified physically plausible variables, including oxygen, flow, thermal diffusion time, fuel pyrolysis temperature, ignition time, and gas or material properties. These findings are diagnostic rather than causal. SHAP decomposes the fitted model; it does not identify the effect of an intervention. Correlated derived variables, uneven missingness, and campaign confounding prevent causal interpretation.

Accordingly, the practical contribution is a reproducible decision-support framework. It freezes splits, isolates tuning within outer training data, verifies paper disjointness, records data fingerprints and feature manifests, reports repeated-fold and paper-level uncertainty, and distinguishes unbiased evaluation from full-data refitting. Within a documented applicability domain, its predictions can rank candidate tests, identify poorly covered regimes, and focus expert attention. They must be accompanied by uncertainty, out-of-distribution checks, threshold rationale, and explicit abstention.

The work therefore establishes a useful but bounded capability: literature-scale machine learning can extract transferable signal for spacecraft fire-safety screening, yet database heterogeneity and source shift remain the dominant obstacles. The path to stronger evidence is not a larger headline score on mixed rows. It is prospective external validation, targeted partial-gravity experiments, a harmonised ontology, uncertainty-aware abstention, and physically constrained models evaluated against genuinely new sources. Until those steps are completed, the system should inform experiments and engineering judgement, never replace them.
